\chapter{Remaining-Time Predictions in Days}
\label{app:lognormal}
The regressor of Section~\ref{subsec:ev:off:models} works on a logarithmic
scale, while the remaining time that reaches the Pareto front is in days.
This appendix goes through the conversion from one to the other step by step.
Throughout, $T$ is the remaining time of a case in days; the values in
seconds used by the implementation are $86\,400\,T$.

\section{What the Regressor Predicts}
The regressor is trained on
\begin{equation}
  X = \log(1 + T)
  \label{eq:app:x}
\end{equation}
instead of on $T$. Remaining times are skewed: most cases end within a few
days, a few take months. On the scale of $X$ the long tail is compressed:
an error of one day on a case that ends in ten days weighs about as much as
an error of ten days on a case that ends in a hundred, and the
distribution is close enough to a Normal for \texttt{RMSEWithUncertainty} to
fit it. The added one keeps $X$ defined at $T = 0$ and makes the scale almost
linear below one day.

For a given input the regressor returns the mean $\mu$ and the variance
$\sigma^2$ of a Normal distribution over $X$. Here $\mu$ is the mean of the
predictions of the members of the virtual ensemble and $\sigma^2$ their total
variance (Equation~\eqref{eq:bg:unc:gaussvar}), multiplied by $s^2$
(Equation~\eqref{eq:bg:unc:sigmafit}).

\section{Going Back to Days}
Inverting~\eqref{eq:app:x} gives $T = e^X - 1$, an increasing function of
$X$. An increasing function maps the median of $X$ onto the median of $T$,
so the median of the remaining time is simply $e^{\mu} - 1$. The mean does
not carry over in the same way. When $X$ is Normal, $e^X$ follows a
log-normal distribution, whose mean is $e^{\mu + \sigma^2/2}$, so
\begin{equation}
  \mathbb{E}[T] = e^{\mu + \sigma^2 / 2} - 1 .
  \label{eq:app:plainmean}
\end{equation}
The point prediction of the method is the mean, since the front compares
the remaining times that the cases are expected to take. The extra term
$\sigma^2/2$ is what makes the mean larger than the median: the long tail of
the log-normal pulls it upwards, and the more uncertain the prediction, the
longer the tail.

This term is also the problem. Table~\ref{tab:app:example} takes a
prediction with $\mu = 2$, whose median is $6.4$ days, and three values of
$\sigma^2$. With $\sigma^2 = 0.25$ the mean is $7.4$ days, close to the
median. With $\sigma^2 = 4$ it is already $53.6$ days, and with $\sigma^2 = 9$
it is $664$ days, almost two years, for a case that most likely ends within
a week. A handful of inputs with a very large predicted variance is enough to
produce such values, and they would end up on the front.

\begin{table}[h]
  \centering
  \small
  \begin{tabular}{rrrr}
    \toprule
    $\sigma^2$ & Median & Mean, Eq.~\eqref{eq:app:plainmean}
      & Mean, truncated at $U = 200$ days \\
    \midrule
    0.25 & 6.4 &   7.4 &  7.4 \\
    4    & 6.4 &  53.6 & 19.9 \\
    9    & 6.4 & 664.1 & 21.2 \\
    \bottomrule
  \end{tabular}
  \caption{Remaining time in days for a prediction with $\mu = 2$ and three
    values of the predicted variance $\sigma^2$.}
  \label{tab:app:example}
\end{table}

\section{Truncating at the Longest Remaining Time}
No case of the training log $\mathcal{L}$ has a remaining time longer than
its maximum $U$, so the predictive distribution is restricted to
$T \le U$, that is, to $X \le b$ with $b = \log(1 + U)$. Truncating the
Normal at $b$ means discarding its values above $b$ and rescaling the
probability left below $b$ so that it sums to one again. Writing $\Phi$ for
the cumulative distribution function of the standard Normal, the
probability of the part that is kept is
\begin{equation}
  P(X \le b) = \Phi(z_b), \qquad z_b = \frac{b - \mu}{\sigma},
  \label{eq:app:zb}
\end{equation}
where $z_b$ is the distance of the bound from the mean, in standard
deviations. The truncated distribution has the density of the Normal, cut at
$b$ and divided by $\Phi(z_b)$:
\begin{equation}
  f_b(x) =
  \begin{cases}
    \dfrac{1}{\Phi(z_b)}\, \dfrac{1}{\sigma\sqrt{2\pi}}\,
    e^{-(x - \mu)^2 / (2\sigma^2)} & \text{if } x \le b,\\[2ex]
    0 & \text{if } x > b.
  \end{cases}
  \label{eq:app:density}
\end{equation}
When the bound lies far above the mean, $z_b$ is large, $\Phi(z_b)$ is
close to one and the truncation changes almost nothing.

\section{Mean of the Truncated Distribution}
The point prediction is $\mathbb{E}[T \mid X \le b] = \mathbb{E}[e^X \mid X
\le b] - 1$, so what is needed is the mean of $e^X$ under the truncated
distribution. The derivation below is done for $e^{nX}$, since the case
$n = 2$ gives the standard deviation as well, and it takes four steps.

\paragraph{Step 1: definition.}
A conditional expectation is an average over the part of the distribution
that satisfies the condition:
\begin{equation}
  \mathbb{E}\bigl[ e^{nX} \mid X \le b \bigr]
  = \frac{\mathbb{E}\bigl[ e^{nX} \cdot \mathbb{I}\{X \le b\} \bigr]}{P(X \le b)} .
  \label{eq:app:step1}
\end{equation}
The indicator $\mathbb{I}\{X \le b\}$ is one when $X \le b$ and zero
otherwise, so it removes the values above $b$ from the expectation in the
numerator, and the division by $P(X \le b)$ turns what is left into an
average over the part that is kept.

\paragraph{Step 2: the integral.}
The expectation in the numerator is the integral of $e^{nx}$ weighted by the
Normal density, and the indicator only moves the upper limit of the integral
from $+\infty$ to $b$. With $P(X \le b) = \Phi(z_b)$,
\begin{equation}
  \mathbb{E}\bigl[ e^{nX} \mid X \le b \bigr]
  = \frac{1}{\Phi(z_b)} \int_{-\infty}^{b} e^{nx}\,
    \frac{1}{\sigma\sqrt{2\pi}}\, e^{-(x - \mu)^2 / (2\sigma^2)}\, dx ,
  \label{eq:app:step2}
\end{equation}
which is the integral of $e^{nx}$ against the truncated density $f_b$ of
Equation~\eqref{eq:app:density}.

\paragraph{Step 3: completing the square.}
The two exponentials in the integral combine into one, whose exponent can be
rewritten as
\begin{equation}
  n x - \frac{(x - \mu)^2}{2\sigma^2}
  = \Bigl( n\mu + \frac{n^2\sigma^2}{2} \Bigr)
    - \frac{\bigl( x - (\mu + n\sigma^2) \bigr)^2}{2\sigma^2} .
  \label{eq:app:square}
\end{equation}
Expanding both sides shows that they are equal.

\paragraph{Step 4: the integral that remains.}
The term in parentheses does not depend on $x$, so it leaves the integral as
the factor $e^{n\mu + n^2\sigma^2/2}$. The rest of the integrand is again the
density of a Normal, with the same variance $\sigma^2$ but with its mean
shifted from $\mu$ to $\mu + n\sigma^2$. Integrated from $-\infty$ to $b$,
it is the probability that this shifted Normal falls below $b$,
\begin{equation}
  \int_{-\infty}^{b} \frac{1}{\sigma\sqrt{2\pi}}\,
  e^{-(x - \mu - n\sigma^2)^2 / (2\sigma^2)}\, dx
  = \Phi\Bigl( \frac{b - \mu - n\sigma^2}{\sigma} \Bigr)
  = \Phi(z_b - n\sigma) .
  \label{eq:app:shifted}
\end{equation}
Putting the factor, the remaining integral and the division by
$\Phi(z_b)$ together,
\begin{equation}
  \mathbb{E}\bigl[ e^{nX} \mid X \le b \bigr]
  = e^{\,n\mu + n^2\sigma^2/2}\; \frac{\Phi(z_b - n\sigma)}{\Phi(z_b)} .
  \label{eq:app:moments}
\end{equation}

\paragraph{The point prediction.}
For $n = 1$,
\begin{equation}
  \mathbb{E}[T \mid T \le U]
  = e^{\mu + \sigma^2/2}\; \frac{\Phi(z_b - \sigma)}{\Phi(z_b)} - 1 .
  \label{eq:app:mean}
\end{equation}
Without the truncation the bound $b$ goes to infinity, both values of $\Phi$
go to one, and Equation~\eqref{eq:app:mean} becomes the log-normal mean of
Equation~\eqref{eq:app:plainmean}. The truncation multiplies that mean by
the correction factor $\Phi(z_b - \sigma) / \Phi(z_b)$. When $\sigma$ is
small the factor is close to one. When $\sigma$ is large, $z_b - \sigma$ is
far in the left tail of the Normal, $\Phi(z_b - \sigma)$ is tiny, and the
factor cancels the growth of $e^{\mu + \sigma^2/2}$. In the example of
Table~\ref{tab:app:example} with $\mu = 2$, $\sigma = 3$ and $U = 200$ days,
$b = \log 201 \approx 5.30$ and $z_b \approx 1.10$, so
$\Phi(z_b) \approx 0.86$ and $\Phi(z_b - \sigma) = \Phi(-1.90) \approx
0.029$. The mean of $e^X$ is then $e^{6.5} \cdot 0.029 / 0.86 \approx 22.2$
instead of $e^{6.5} \approx 665$, and the remaining time $21.2$ days instead
of $664$.

The result is finally clipped to $[0, U]$. For a large $\sigma$ the factor
$e^{\mu + \sigma^2/2}$ can overflow while $\Phi(z_b - \sigma)$ underflows,
so the implementation computes the logarithm of~\eqref{eq:app:mean} and
exponentiates only at the end.

\section{Standard Deviation}
The variance of $T$ equals the variance of $e^X$, since the two differ only
by a constant. With~\eqref{eq:app:moments} for $n = 1$ and $n = 2$,
\begin{equation}
  \mathrm{sd}(T \mid T \le U)
  = \sqrt{ \mathbb{E}\bigl[ e^{2X} \mid X \le b \bigr]
         - \mathbb{E}\bigl[ e^{X} \mid X \le b \bigr]^2 } .
  \label{eq:app:sd}
\end{equation}

\section{Quantiles and Intervals}
Integrating the truncated density~\eqref{eq:app:density} from $-\infty$ up to
a value $x \le b$ gives its cumulative distribution function,
\begin{equation}
  P(X \le x \mid X \le b) = \frac{\Phi\bigl( (x - \mu)/\sigma \bigr)}{\Phi(z_b)},
  \qquad x \le b .
  \label{eq:app:cdf}
\end{equation}
Setting it equal to $q$ and solving for $x$ gives the $q$-quantile of $X$,
\begin{equation}
  x_q = \mu + \sigma\, \Phi^{-1}\bigl( q\, \Phi(z_b) \bigr) .
  \label{eq:app:quantile}
\end{equation}
Since $T$ increases with $X$, the $q$-quantile of the remaining time is
$e^{x_q} - 1$. The median of the truncated distribution is the case
$q = 0.5$, and an interval that leaves a probability $\alpha/2$ on each side
goes from the quantile $q = \alpha/2$ to the quantile $q = 1 - \alpha/2$.

\section{Comparisons in the Log Space}
The confidence filter (Section~\ref{subsec:met:online:prob}) needs the
probability that a candidate $C$ shortens the remaining time with respect to
the baseline $B$. Since $T$ is an increasing function of $X$, $T_C < T_B$
holds exactly when $X_C < X_B$, so the probability is the same on the two
scales. The filter therefore compares the members directly on the scale of
$X$, where their predictions are Normal, and needs none of the conversions
above.

\section{Why This Scale and the Truncation}
\label{sec:app:choice}
The scale of~\eqref{eq:app:x} was chosen by training the regressor of
Section~\ref{subsec:ev:off:models} three times on each log, with the same
search space and number of trials, on three scales: the remaining time in
seconds as it is, $\log(1 + \text{seconds})$, and $\log(1 + T)$ with $T$ in
days. The predictions are always converted to days, and the models are
compared on the prefixes of the cases in $\mathcal{L}^{\mathrm{test}}$.
Table~\ref{tab:app:scales} reports the errors of the point prediction and
the calibration of the recalibrated models, computed on each model's own
scale.

\begin{table}[h]
  \centering
  \small
  \begin{tabular}{llrrrrr}
    \toprule
    Log & Scale & MAE & RMSE & $\pm\sigma$ & $\pm 2\sigma$ & ENCE \\
        &       & (days) & (days) & (\%) & (\%) & (\%) \\
    \midrule
    BAC          & seconds            &  5.65 &   18.23 & 96.2 & 98.5 & 36.0 \\
                 & $\log(1 + \text{seconds})$ &  5.61 &   44.94 & 84.0 & 97.9 & 21.9 \\
                 & $\log(1 + \text{days})$    &  5.10 &   18.05 & 87.8 & 97.8 & 28.4 \\
    \midrule
    BPIC2012     & seconds            &  7.82 &   11.96 & 74.8 & 93.8 & 28.4 \\
                 & $\log(1 + \text{seconds})$ &  8.47 &   15.47 & 67.4 & 97.1 & 14.9 \\
                 & $\log(1 + \text{days})$    &  7.86 &   11.94 & 63.7 & 95.8 & 14.1 \\
    \midrule
    bpi17-before & seconds            &  9.32 &   13.80 & 76.4 & 91.9 & 32.1 \\
                 & $\log(1 + \text{seconds})$ &  9.11 &   13.20 & 80.2 & 98.8 & 19.1 \\
                 & $\log(1 + \text{days})$    & 15.91 &  684.32 & 58.7 & 95.7 & 13.7 \\
    \midrule
    bpi17-after  & seconds            &  8.77 &   13.38 & 72.4 & 90.4 & 43.3 \\
                 & $\log(1 + \text{seconds})$ & 22.86 & 1650.86 & 84.1 & 98.5 & 25.0 \\
                 & $\log(1 + \text{days})$    &  8.31 &   12.79 & 74.1 & 97.6 & 17.3 \\
    \bottomrule
  \end{tabular}
  \caption{The regressor trained on three scales. MAE and RMSE of the point
    prediction, in days, with the mean of the log-normal of
    Equation~\eqref{eq:app:plainmean} for the two logarithmic scales; coverage
    of the $\pm\sigma$ and $\pm 2\sigma$ intervals, whose targets are 68.3\%
    and 95.4\%; ENCE. Calibration measures after the scaling by $s$.}
  \label{tab:app:scales}
\end{table}

On the scale of seconds the point errors are acceptable, but the
uncertainty is not: the ENCE is the highest on every log, between 28\% and
43\%, and in BAC the $\pm\sigma$ interval holds 96\% of the cases instead of
68\%, so the predicted standard deviations are far too wide. A Normal over
seconds also gives a positive probability to negative remaining times. The
two logarithmic scales calibrate better, with an ENCE between 14\% and 28\%,
but their point errors show a few enormous values: an RMSE of 684 days on
bpi17-before for $\log(1 + \text{days})$ and of 1\,651 days on bpi17-after
for $\log(1 + \text{seconds})$, against MAEs of 16 and 23 days. The cases
behind these values are very few, 19 of the 18\,085 test prefixes of
bpi17-before and 31 of the 26\,095 of bpi17-after, about one in a thousand,
but their predictions are absurd: the largest reaches 89\,029 days on
bpi17-before, more than two centuries. They are the predictions with a very
large variance of Table~\ref{tab:app:example}, and a handful of them is
enough to raise the mean error, and above all the RMSE, by orders of
magnitude, while the median of the same predictions stays reasonable.

Between the two logarithmic scales, the one in days was chosen. Both suffer
from a few extreme predictions, on different logs, so the errors of the point
prediction do not separate them, while the scale in days has the lower ENCE
on three logs out of four. Moreover, on the scale in seconds a difference of
a few minutes between two very short remaining times weighs as much as a
difference of days between long ones, while the scale in days is almost
linear below one day.

The truncation at $U$ then removes the extreme predictions of the chosen
scale. Table~\ref{tab:app:trunc} shows its effect on the same models,
without retraining them. It brings the 19 predictions of bpi17-before above
$U$ back below it, the largest falling from 89\,029 to 40 days and the RMSE
from 684 to 14 days, and it leaves the other three logs almost unchanged,
since none of their predictions exceeded $U$.

\begin{table}[h]
  \centering
  \small
  \setlength{\tabcolsep}{4pt}
  \begin{tabular}{lrrrrrrrr}
    \toprule
    Log & $U$ & \multicolumn{2}{c}{MAE (days)} & \multicolumn{2}{c}{RMSE (days)}
        & Above & \multicolumn{2}{c}{Largest (days)} \\
    \cmidrule(lr){3-4} \cmidrule(lr){5-6} \cmidrule(lr){8-9}
        & (days) & plain & trunc. & plain & trunc. & $U$ & plain & trunc. \\
    \midrule
    BAC          & 503.6 &  5.10 & 5.09 &  18.05 & 18.05 &  0 &      38.7 & 33.0 \\
    BPIC2012     &  91.4 &  7.86 & 7.91 &  11.94 & 12.03 &  0 &      30.6 & 26.5 \\
    bpi17-before & 135.5 & 15.91 & 9.15 & 684.32 & 13.67 & 19 & 89\,029.2 & 40.2 \\
    bpi17-after  & 101.0 &  8.31 & 8.42 &  12.79 & 12.96 &  0 &      33.5 & 33.5 \\
    \bottomrule
  \end{tabular}
  \caption{Effect of the truncation at the longest training remaining time
    $U$ on the models trained on $\log(1 + \text{days})$ of
    Table~\ref{tab:app:scales}. MAE and RMSE without (plain) and with
    (trunc.) the truncation, number of test predictions above $U$ without
    it, and largest test prediction without and with it.}
  \label{tab:app:trunc}
\end{table}
